\chapter{Derivation of the Adjoint System}
\label{app:adjoint}
\renewcommand{\arraystretch}{1.0}

In this Appendix we derive the adjoint system (\ref{eq:euler:voigt:adjoint}) using the duality-pairing relation (\ref{eq:duality:pairing}). We drop the subscript $\alpha$ from $\ua$, $\usa$, and $\upa$ to simplify the notation. {\color{red} Additionally, we use a simplified integral notation where $\Theta = (\Omega \times [0,T])$, i.e.
$$\int_0^T\int_{\Omega} \u \dxdt = \int_\Theta \u \dxdt. $$
}We begin by expanding (\ref{eq:duality:pairing}) into individual integrals which will be evaluated separately

\begin{subequations}
\begin{align}
\left(\L\begin{bmatrix}
\u'\\
p'
\end{bmatrix}, \begin{bmatrix}
\u^*\\
p^*
\end{bmatrix} \right) =& \int_{\Theta} \partial_t \u' \cdot \u^* \dxdt \label{eq:integral:1} \\
+& \int_{\Theta} (I\text{-}\,\alpha^2\Delta)^{-1}(\u\cdot\bn\u')\cdot\u^* \dxdt \label{eq:integral:2} \\
+& \int_{\Theta} (I\text{-}\,\alpha^2\Delta)^{-1}(\u'\cdot\bn\u)\cdot\u^* \dxdt \label{eq:integral:3} \\
+& \int_{\Theta} \bn p' \cdot \u^* \dxdt \label{eq:integral:4} \\ 
+& \int_{\Theta} (\bn \cdot \u')p^* \dxdt. \label{eq:integral:5}
\end{align}
\end{subequations}

\noindent Integrating (\ref{eq:integral:1}) by part with respect to time gives
$$\int_{\Theta} \partial_t \u' \cdot \u^* \,d\x\, dt = \int_{\Omega}\u^*\cdot \u'|_{t=0}^{t=T} d\x - \int_{\Theta} \partial_t \u^* \cdot \u' \dxdt.$$
\\~\\
\noindent To evaluate (\ref{eq:integral:2}) we first split the integrand into its components, and use Einstein's notation to simplify.

$$\int_{\Theta}(I\text{-}\,\alpha^2\Delta)^{-1}(\u\cdot\bn\u')\cdot \u^* \dxdt = \int_{\Theta}(I\text{-}\,\alpha^2\Delta)^{-1}u_i\partial_i u'_j u_j^* \dxdt.$$

\noindent Integrating (\ref{eq:integral:2}) by part with respect to space gives
\begin{subequations}
\begin{align*}
\int_{\Theta} (I\text{-}\,\alpha^2\Delta)^{-1}(\u\cdot\bn\u')\cdot\u^* \dxdt = &\int_0^T (I\text{-}\,\alpha^2\Delta)^{-1} u_i u_j^* u'_j |_{\partial \Omega} \,dt - \\ &\int_{\Theta}  u'_j \partial_i((I\text{-}\,\alpha^2\Delta)^{-1} u_i u_j^*) \dxdt.
\end{align*}
\end{subequations}
\noindent Because $\Omega$ has periodic boundaries, the first integral is zero, and we can expand the second integral.

$$\int_{\Theta} (I\text{-}\,\alpha^2\Delta)^{-1}(\u\cdot\bn\u')\cdot\u^* \dxdt = \int_{\Theta} u'_j\ls u_i\partial_i((I\text{-}\,\alpha^2\Delta)^{-1} u_i) - (I\text{-}\,\alpha^2\Delta)^{-1} u_j^*\partial_i u_i \rs \dxdt.$$

\noindent The flow is divergence free and so $\partial_i u_i = 0$. We can convert the second integral back to vector notation to obtain

$$\int_{\Theta}(I\text{-}\,\alpha^2\Delta)^{-1}(\u\cdot\bn\u')\cdot \u^* \dxdt = -\int_{\Theta} (\u\cdot\bn(I\text{-}\,\alpha^2\Delta)^{-1}\u^*)\cdot \u' \dxdt.$$
\\~\\
\noindent To evaluate (\ref{eq:integral:3}) we first split the integrand into its components, and use Einstein's notation to simplify.

$$\int_{\Theta}(I\text{-}\,\alpha^2\Delta)^{-1}(\u\cdot\bn\u')\cdot \u^* \dxdt = \int_{\Theta}(I\text{-}\,\alpha^2\Delta)^{-1}u'_i(\partial_i u_j) u_j^* \dxdt.$$

\noindent The flow is divergence free and so we can convert express this as

$$\int_{\Theta}(I\text{-}\,\alpha^2\Delta)^{-1}(\u\cdot\bn\u')\cdot \u^* \dxdt = \int_{\Theta}(I\text{-}\,\alpha^2\Delta)^{-1}\partial_i(u'_i u_j) u_j^* \dxdt.$$

\noindent Integrating (\ref{eq:integral:3}) by part with respect to space gives.
\begin{subequations}
\begin{align*}
\int_{\Theta}(I\text{-}\,\alpha^2\Delta)^{-1}(\u\cdot\bn\u')\cdot \u^* \dxdt = &\int_0^T (I\text{-}\,\alpha^2\Delta)^{-1} u_j^* u'_i u_j \,dt - \\ 
&\int_{\Theta} u'_i u_j \partial_i(I\text{-}\,\alpha^2\Delta)^{-1} u_j^* \dxdt
\end{align*}
\end{subequations}

\noindent Because $\Omega$ has periodic boundaries, the first integral is zero. We can convert the second integral back to vector notation to obtain

$$\int_{\Theta}(I\text{-}\,\alpha^2\Delta)^{-1}(\u\cdot\bn\u')\cdot \u^* \dxdt = -\int_{\Theta} (\bn(I\text{-}\,\alpha^2\Delta)^{-1}\u^*)^\intercal \u \cdot \u' \dxdt.$$
\\~\\
\noindent To evaluate (\ref{eq:integral:4}) we first split the integrand into its components, and use Einstein's notation to simplify.

$$\int_{\Theta} \bn p' \cdot \u^* \dxdt = \int_{\Theta} \partial_i p'\u_i^* \dxdt.$$

\noindent Integrating (\ref{eq:integral:4}) by part with respect to space gives

$$\int_{\Theta} \bn p' \cdot \u^* \,d\x\,dt = \int_0^T \u^* p' |_{\partial \Omega} \,dt - \int_{\Theta} p' \partial_i u_i^* \dxdt.$$

\noindent Because $\Omega$ has periodic boundaries, the first integral is zero, and we can convert the second integral back to vector notation to obtain

$$\int_{\Theta} \bn p' \cdot \u^* \,d\x\,dt = -\int_{\Theta} p' (\bn\cdot \u^*) \dxdt.$$
\\~\\
\noindent To evaluate (\ref{eq:integral:5}) we first split the integrand into its components, and use Einstein's notation to simplify.

$$\int_{\Theta} (\bn \cdot \u')p^* \,d\x\,dt = \int_{\Theta} (\partial_i \u'_i)p^* \dxdt$$

\noindent Integrating (\ref{eq:integral:5}) by part with respect to space

$$\int_{\Theta} (\bn \cdot \u')p^* \,d\x\,dt = \int_0^T p^*\u'|_{\partial \Omega} \,dt - \int_{\Theta} \u'_i \partial_i p^* \dxdt.$$

\noindent Because $\Omega$ has periodic boundaries, the first integral is zero, and we can convert the second integral back to vector notation to obtain

$$\int_{\Theta} (\bn \cdot \u')p^* \,d\x\,dt = -\int_{\Theta} \u' \cdot \bn p^* \dxdt.$$

\noindent Combining the integrals, we can express the LHS of (\ref{eq:duality:pairing}) as

\begin{subequations}
\begin{align}
\left(\L\begin{bmatrix}
\u'\\
p'
\end{bmatrix}, \begin{bmatrix}
\u^*\\
p^*
\end{bmatrix} \right) =& \int_{\Omega}\u^* \cdot \u'|_{t=0}^{t=T} d\x - \int_{\Theta} \u'\cdot \partial_t \u^* \dxdt. \\
-& \int_{\Theta} (\bn(I\text{-}\,\alpha^2\Delta)^{-1}\u^*\cdot \u)\cdot \u' \dxdt. \\
-& \int_{\Theta} (\bn(I\text{-}\,\alpha^2\Delta)^{-1}\u^*)^\intercal \u \cdot \u' \dxdt \\
-& \int_{\Theta} p' (\bn\cdot \u^*) \dxdt \\ 
-& \int_{\Theta} \u' \cdot \bn p^* \dxdt.
\end{align}
\end{subequations}

\noindent Using this and the right-hand side (RHS) of (\ref{eq:duality:pairing}) we obtain the adjoint system,
$$\L^*\begin{bmatrix}
\usa\\
\psa
\end{bmatrix} := \begin{bmatrix}
-\partial_t\usa-(\bn (I-\alpha^2\Delta)^{-1}\usa + (\bn (I-\alpha^2\Delta)^{-1}\usa)^T)\ua-\bn \psa \\
-\bn \cdot \usa
\end{bmatrix} = \begin{bmatrix}
0\\
0
\end{bmatrix}$$

\noindent To derive the terminal condition of the adjoint system we use the RHS of (\ref{eq:duality:pairing}).
$$\int_\Omega \usa(T) \cdot \upa(T) \,dx = \int_\Omega \usa(0) \cdot \upa(0) \,d\x.$$
From the two forms of the G\^{a}teaux differential, (\ref{eq:Gateaux:differential:inner:product}) and (\ref{eq:Gateaux:differential:inner:product:2}), we obtain 
$$\int_\Omega \usa(T) \cdot \upa(T) \,dx = 2\int_\Omega |D|\usa(T) \cdot |D|\upa(T) \,d\x.$$
Therefore,
$$\usa(T) = 2|D|^2\ua(T).$$
